(working in @@M0@@ by means of the identifications 1.4).

In view of 1.6.ii), @@M1@@, and the @@M2@@ define an element @@M3@@ of @@M4@@, since this is a separated complete group; we therefore put @@M5@@.

Finally, let @@M6@@, and @@M7@@; we put @@M8@@, where @@M9@@ is the homomorphism from @@M10@@ to @@M11@@ defined in 1.2. Since we have already defined @@M12@@ on the elements @@M13@@, we must verify that the two definitions give the same result, that is:
@@M14@@
Now, for every @@M15@@, we have: @@M16@@. Hence, for every @@M17@@; @@M18@@. Using the identifications of 1.4 and the gluing of the @@M19@@ for varying @@M20@@, we obtain: @@M21@@.
Thus, finally:
@@M22@@
Since @@M23@@ is separated, we obtain the required equality.

We have thus defined a system of maps @@M24@@; to verify that they constitute a morphism of functors (with values in @@M25@@), we must show that for every morphism @@M26@@, the diagram:
@@M27@@
